\documentclass[instructions.tex]{subfiles}
\begin{document}
 
\chapter{De l’orgueil.}
 
L’ORGUEIL est un désir d’être estimé, d’être élevé au-dessus des autres, de dominer, de vivre dans l’indépendance. Ce vice est, 1.o haïssable devant Dieu et devant les hommes; 2.o d’autant plus haïssable, qu’il n’est fondé sur aucune raison qui puisse le justifier\footnote{Odibilis coràm Deo est et hominibus superbia. Eccli. 10. 7.}.
 
\primo Dieu ne peut souffrir l’orgueilleux. A peine Lucifer veut-il s’élever, qu’il est précipité au fond des abîmes. A mesure qu’un esprit superbe s’enfle et s’élève, Dieu lui résiste\footnote{Superbis resistit.}; il lui retire sa grâce, et l’abandonne à lui-même. C’est pour cela qu’il livra les philosophes du paganisme à leur sens réprouvé, à des passions d’ignominie, et qu’il abandonne aujourd’hui les esprits forts de ce siècle, enflés de leurs talens, aux désirs de leur cœur, à un aveuglement qui ébranle leur croyance, qui détruit en eux les sentimens de la religion et de la foi, selon cette parole du Sauveur : Comment pouvez-vous croire, vous qui aimez à recevoir de la gloire les uns des autres\footnote{Quomodò vos potestis credere, qui gloriam ab invicem accipitis? Joan. 5. 44.}?
 
Ce vice est haï des hommes, parce qu’un orgueilleux met le trouble partout où il se trouve. Un esprit superbe a mis la division dans le ciel parmi les anges. Il ne faut qu’un savant orgueilleux pour mettre le schisme dans l’Eglise : il ne faut aussi qu’un esprit dominant, un génie superbe et hautain dans une famille, dans une communauté, dans une ville, pour y mettre le trouble et le désordre. Il n’y a ni confiance, ni amitié sincère, ni paix, ni concorde parmi les orgueilleux\footnote{Inter superbos semper jurgia sunt. Prov. 13. 10.}, parce que l’orgueilleux veut l’emporter sur tous. Idolâtre de lui-même, il n’estime que ses sentimens : d’où il arrive que personne ne peut le souffrir, comme il ne peut souffrir personne. Si un homme vain savait ce qu’on pense de lui, rien ne serait plus propre à rabattre son orgueil.
 
\secundo Ce vice est d’autant plus détestable, qu’il n’est fondé sur rien qui puisse l’excuser. On s’élève sottement, quand on n’a rien qui relève, et rien de bon qui soit de nous. Qu’avez-vous de vous-même que la misère et le néant? L’homme n’a de son fond, dit un concile, que l’erreur et le péché. Y a-t-il là de quoi vous enfler d’orgueil? Tout ce que vous avez de bon, la santé, l’esprit, les talens, ne sont-ce pas des faveurs dont Dieu vous a gratuitement pourvu?
 
La vertu même et vos bonnes œuvres sont tellement un effet de la bonté de Dieu, que ce sont ses propres dons qu’il couronne, dit saint Augustin, quand il couronne vos mérites. La moindre bonne œuvre dépend plus de la grâce et du concours de Dieu, que les rayons ne dépendent du soleil. Les rayons pourraient, par miracle, subsister sans le soleil, mais il n’est point de créature qui puisse quelque chose sans le secours de Dieu.
 
C’est donc faire injure à Dieu, que de vous approprier la gloire du bien que vous faites; c’est lui qui en est le premier auteur. Remerciez-le, réjouissez-vous même de ce qu’il se sert de vous pour procurer sa gloire; mais aussi humiliez-vous de ce qu’il se sert d’un instrument aussi misérable que vous. En faisant de grandes choses pour Dieu, il ne devient pas votre redevable, vous ne faites que votre devoir. Ce que vous faites n’ajoute rien à sa félicité, puisqu’il n’a pas besoin de vous.
 
Tout ce que nous devons, après avoir fait nos efforts pour lui plaire, c’est d’avouer humblement que nous sommes des serviteurs inutiles; que nous n’avons fait que ce que nous devions faire; que nos œuvres les plus saintes sont, à ses yeux, comme un linge souillé et immonde; que nous sommes même indignes qu’il nous souffre à son service. Nous aurons du moins cette consolation, que si nous ne pouvons rien faire à Dieu pour le rendre plus heureux, nous pourrons du moins le glorifier par notre humilité.
 
Par rapport aux hommes, vous n’avez de même aucun sujet de vous préférer à personne. Si vous êtes plus élevé, si vous avez des inférieurs ou des sujets, vous ne devez pas vous estimer plus qu’eux : si c’eût été la volonté de Dieu, votre domestique ou votre sujet aurait été votre maître, et vous auriez été son domestique ou son sujet.
 
Eussiez-vous même de grandes vertus, vous ne devez pas vous estimer plus que les scélérats. Ne savez-vous pas qu’un grand pécheur peut devenir un grand saint; que tel que vous méprisez pour ses désordres, sera peut-être un jour dans le ciel, et que vous-même, fussiez-vous l’homme le plus vertueux, vous pouvez devenir, en un moment, un réprouvé? Si vous avez quelque chose de bon, dit saint Paul, ne l’avez-vous pas reçu de Dieu? pourquoi vous en glorifiez-vous?
 
Si ce que nous sommes, et ce qui est bon dans nous, ne doit pas nous donner de l’orgueil, ce qui est au-dehors doit encore moins nous en inspirer. Tirer vanité de ses habits, de sa beauté, de ses richesses, de son crédit, de son rang, c’est se méconnaître, et oublier ce que l’on est. Un animal chargé d’ornemens magnifiques, monté par un César, n’est toujours qu’une bête comme les autres animaux de son espèce. Une pierre richement taillée, posée sur le frontispice d’un palais, n’est toujours qu’une pierre de la nature de celles qui sont dans la carrière. De même, un homme environné de magnificence et d’éclat, élevé par sa fortune et ses emplois, n’est toujours qu’un homme comme les autres; avec cette différence, qu’il est peut-être plus grand pécheur, et que ceux qu’il méprise sont souvent plus estimés, plus dignes que lui de son rang et de sa fortune.
 
Les nobles, les personnes de qualité, doivent se ressouvenir que, devant Dieu, ils ne sont pas plus que les autres mortels. Formés de boue, tirés de la même poussière, leur origine les rend tous égaux. Qu’ils prennent garde que leur élévation ne serve un jour à les confondre, et à leur attirer un plus sévère châtiment, selon cette menace de Dieu : Les puissans seront tourmentés puissamment\footnote{Potentes potenter tormenta patientur. Sap. 6. 7.}. Devant les hommes, leur noblesse ne sert qu’à les rendre méprisables, si, au lieu de la soutenir par la vertu, ils la soutiennent par la fierté et par le vice. Autant la vertu donne d’éclat à leur naissance, autant le vice la dégrade. C’est la vertu qui rend l’homme respectable et vraiment noble\footnote{Moribus et vitâ nobilitatur homo.}.
 
Finissons cet important chapitre par ces belles paroles de Tobie : Prenez garde, mon fils, que l’orgueil ne domine jamais dans votre cœur et dans vos paroles, car il est le commencement de toute perdition\footnote{Tob. 4. 11.}.
 
\section{Résolutions}
 
\primo Puisque l’orgueil déplaît si fort à Dieu, et qu’il lui ravit la gloire qui lui est due; puisqu’il produit de si mauvais effets dans la société, et qu’il couvre d’un juste ridicule celui qui s’y abandonne, je le combattrai en moi toute ma vie.
 
\secundo Mais afin de le faire avec plus de succès, je me ressouviendrai souvent que tout ce que je possède de bon, je l’ai reçu de la bonté de Dieu, comme auteur de la nature ou de la grâce.
 
\tertio Ainsi, je ne me glorifie-rai plus de rien; rendant, au contraire, gloire au Seigneur de tout le bien que je reconnais en moi. 

\prayer{Mon Dieu, ôtez de mon cœur l’orgueil, la vanité, la présomption, et tous les autres vices qui viennent à la suite de ce monstre.}
 
\end{document}
